
\section{Relacionar documentación y estado físico}
Una revisión de aeronavegabilidad no se reduce a fechas. Identidad, configuración, estado físico y cumplimiento deben reconciliarse. Un historial aparentemente completo puede omitir una instalación; un buen aspecto exterior puede coexistir con una limitación vencida. Una inspección general tampoco sustituye todos los registros.

M.A.901 y M.A.903 de la consolidación de 2026 permiten distinguir condiciones de emisión o recomendación y proceso de revisión; M.A.902 trata validez. Los privilegios y documentos dependen del escenario. No se atribuye la misma facultad a cualquier organización por trabajar habitualmente con la aeronave. ML necesita su propio análisis.
\section{Documentos con preguntas distintas}
Certificado de tipo, CRS y ARC tienen objetos diferentes. La existencia de uno no renueva automáticamente otro. Un CRS tras una tarea no añade por sí solo otro periodo de ARC; un ARC no demuestra correcta ejecución de cualquier mantenimiento posterior. Certificado de aeronavegabilidad y ARC también se distinguen.

Una fecha en planificación no constituye una decisión de autoridad. Una ventana permitida para anticipar revisión se usa en su escenario exacto; no se convierte en tolerancia general de mantenimiento o límites de vida.
\section{Resolver una discrepancia de estado}
Si una hoja dice «cumplido» sin evidencia, se identifica tarea, fuente, revisión, configuración y registro que sostendría el estado. Una inspección visual genérica no equivale a una acción específica de AD. Se registra el hallazgo y se aclara por el proceso pertinente, sin fabricar un cierre con datos supuestos.

«Pendiente de demostrar» puede ser la conclusión correcta aunque la fecha esté bien calculada. Localiza qué prueba falta, sin convertir ausencia de evidencia en prueba automática de una condición física concreta.

\Needspace{.75\textheight}
\section{Vuelos de comprobación y autorizaciones de vuelo}
Un vuelo posterior al mantenimiento puede servir para comprobar funciones que no se evalúan completamente en tierra. La necesidad se define por los datos y el proceso; M.A.301(i) incluye vuelos de comprobación de mantenimiento cuando sean necesarios. Esto no significa que todo mantenimiento obligue a un vuelo, ni que cualquier vuelo de prueba pueda realizarse sin analizar condición de aeronave, operación y tripulación.

Parte 21, subparte P, regula permisos de vuelo. El supuesto de 21.A.701(a) se refiere a aeronaves que no cumplen o no han demostrado cumplir requisitos aplicables, pero son capaces de vuelo seguro bajo condiciones definidas y para los fines enumerados. Entre esos fines figura trasladarse donde se realizará mantenimiento o revisión; también hay un supuesto para localizar averías o comprobar sistemas tras mantenimiento. El propósito por sí solo no demuestra que todos los requisitos de emisión estén satisfechos \cite{EU748}.

El punto 21.A.708 trata condiciones de vuelo, incluida configuración, restricciones, procedimientos, programa de ensayos cuando proceda y disposiciones de aeronavegabilidad continuada. 21.A.710 regula su aprobación y 21.A.711 la emisión del permiso según autoridad o privilegios de organización pertinentes. Aprobar condiciones y emitir permiso son pasos relacionados con responsabilidades propias. No se atribuye esa facultad a cualquier taller por tener Parte 145.

Un permiso no transforma una aeronave en plenamente conforme para toda operación. Conserva finalidad, condiciones y restricciones; cambios de configuración pueden afectar su validez y tratamiento. En un estudio se identifica primero por qué se necesita el vuelo, qué condición tiene la aeronave y qué régimen y datos se aplican. Después se determina la ruta competente, sin improvisar un perfil de ensayo ni usar el libro como autorización.
\caso{M16-B. Traslado académico a un taller}{Una aeronave ficticia tiene una discrepancia y se propone volar al lugar de mantenimiento. Antes de calcular combustible o fecha se pregunta si está en condición de volar bajo el régimen ordinario o si necesita otra vía. La mera conveniencia del traslado no acredita seguridad ni concede permiso. Se reúnen configuración, discrepancia, datos y finalidad para analizar condiciones y aprobación en su marco.}

\repaso{¿Puede una ventana de revisión ampliar una inspección del programa?}{No se deduce tal autorización. Cada disposición tiene objeto y condiciones; se verifica el requisito de la tarea.}
